
\documentclass{article}
\usepackage{colortbl}
\usepackage{makecell}
\usepackage{multirow}
\usepackage{supertabular}

\begin{document}

\newcounter{utterance}

\centering \large Interaction Transcript for game `bbh', experiment `multistep\_arithmetic\_two', episode 213 with reflexion{-}Qwen3.5{-}27B.
\vspace{24pt}

{ \footnotesize  \setcounter{utterance}{1}
\setlength{\tabcolsep}{0pt}
\begin{supertabular}{c@{$\;$}|p{.15\linewidth}@{}p{.15\linewidth}p{.15\linewidth}p{.15\linewidth}p{.15\linewidth}p{.15\linewidth}}
    \# & \multicolumn{2}{c}{Player} && \multicolumn{2}{c}{Game Master} \\
    \hline

    \theutterance \stepcounter{utterance}  
    & & & \multicolumn{4}{p{0.6\linewidth}}{
        \cellcolor[rgb]{0.9,0.9,0.9}{
            \makecell[{{p{\linewidth}}}]{
                \texttt{\tiny{[P1$\langle$GM]}}
                \texttt{Solve multi{-}step arithmetic problems.} \\
\\ 
\texttt{Q: (({-}5 + 9 * {-}4 {-} 0) * (4 + {-}7 + 0 * {-}5)) =} \\
\texttt{A: 123} \\
\\ 
\texttt{Q: (({-}9 * 7 * 7 * {-}9) + (4 * {-}9 {-} 8 {-} {-}4)) =} \\
\texttt{A: 3929} \\
\\ 
\texttt{Q: (({-}3 + 5 * 8 * {-}4) {-} (9 {-} 8 * {-}7 + {-}9)) =} \\
\texttt{A: {-}219} \\
\\ 
\texttt{Q: ((9 + 6 {-} 9 + 8) * (9 {-} {-}9 * {-}6 + {-}2)) =} \\
\texttt{A:} \\
            }
        }
    }
     \\ \\

    \theutterance \stepcounter{utterance}  
    & \multicolumn{4}{p{0.6\linewidth}}{
        \cellcolor[rgb]{0.561,0.737,0.561}{
            \makecell[{{p{\linewidth}}}]{
                \texttt{\tiny{[P1$\rangle$GM]}}
                \texttt{To solve the expression \$((9 + 6 {-} 9 + 8) * (9 {-} {-}9 * {-}6 + {-}2))\$, we will break it down into two parts based on the parentheses and follow the order of operations (PEMDAS/BODMAS).} \\
\\ 
\texttt{**Step 1: Evaluate the first part \$(9 + 6 {-} 9 + 8)\$**} \\
\texttt{This part contains only addition and subtraction, so we calculate from left to right:} \\
\texttt{1. \$9 + 6 = 15\$} \\
\texttt{2. \$15 {-} 9 = 6\$} \\
\texttt{3. \$6 + 8 = 14\$} \\
\\ 
\texttt{So, the first part equals **14**.} \\
\\ 
\texttt{**Step 2: Evaluate the second part \$(9 {-} {-}9 * {-}6 + {-}2)\$**} \\
\texttt{This part contains multiplication, addition, and subtraction. We must perform the multiplication first.} \\
\texttt{1. Calculate the multiplication: \${-}9 * {-}6\$} \\
\texttt{   {-} A negative times a negative is a positive.} \\
\texttt{   {-} \$9 * 6 = 54\$} \\
\texttt{   {-} So, \${-}9 * {-}6 = 54\$.} \\
\\ 
\texttt{Now substitute this back into the expression:} \\
\texttt{\$9 {-} 54 + {-}2\$} \\
\\ 
\texttt{Next, perform addition and subtraction from left to right:} \\
\texttt{2. \$9 {-} 54 = {-}45\$} \\
\texttt{3. \${-}45 + {-}2 = {-}47\$} \\
\\ 
\texttt{So, the second part equals **{-}47**.} \\
\\ 
\texttt{**Step 3: Multiply the results of the two parts**} \\
\texttt{Now we multiply the result from Step 1 by the result from Step 2:} \\
\texttt{\$14 * {-}47\$} \\
\\ 
\texttt{Calculation:} \\
\texttt{\$14 * 47 = 14 * (50 {-} 3) = 700 {-} 42 = 658\$} \\
\texttt{Since one number is positive and the other is negative, the result is negative.} \\
\texttt{\$14 * {-}47 = {-}658\$} \\
\\ 
\texttt{**Final Answer:**} \\
\texttt{{-}658} \\
            }
        }
    }
    & & \\ \\

    \theutterance \stepcounter{utterance}  
    & & & \multicolumn{2}{p{0.3\linewidth}}{
        \cellcolor[rgb]{0.9,0.9,0.9}{
            \makecell[{{p{\linewidth}}}]{
                \texttt{\tiny{[GM$|$GM]}}
                \texttt{Target: {-}658} \\
            }
        }
    }
    & & \\ \\

    \theutterance \stepcounter{utterance}  
    & & & \multicolumn{2}{p{0.3\linewidth}}{
        \cellcolor[rgb]{0.9,0.9,0.9}{
            \makecell[{{p{\linewidth}}}]{
                \texttt{\tiny{[GM$|$GM]}}
                \texttt{game\_result = LOSE} \\
            }
        }
    }
    & & \\ \\

    \theutterance \stepcounter{utterance}  
    & & \multicolumn{4}{p{0.6\linewidth}}{
        \cellcolor[rgb]{0.95,0.95,0.95}{
            \makecell[{{p{\linewidth}}}]{
                \texttt{\tiny{[P1$\leftrightarrow$P]}}
                \texttt{Post{-}episode Reflexion (not sent to GM):} \\
\texttt{New plan: My previous attempt failed because I misinterpreted the expression `9 {-} {-}9 * {-}6`. I calculated `{-}9 * {-}6` as `54` (positive), but the expression actually contains three negative signs in that sequence: the subtraction operator, the first negative number, and the second negative number. The term is `9 {-} (({-}9) * ({-}6))`, which is `9 {-} 54 = {-}45`. Wait, let me re{-}evaluate the specific error. The expression is `9 {-} {-}9 * {-}6 + {-}2`.} \\
\texttt{Order of operations: Multiplication first. `{-}9 * {-}6 = 54`.} \\
\texttt{Then subtraction/addition: `9 {-} 54 + {-}2`.} \\
\texttt{`9 {-} 54 = {-}45`.} \\
\texttt{`{-}45 + {-}2 = {-}47`.} \\
\texttt{My calculation of the second part ({-}47) was actually correct.} \\
\texttt{Let's re{-}check the first part: `9 + 6 {-} 9 + 8`.} \\
\texttt{`9 + 6 = 15`.} \\
\texttt{`15 {-} 9 = 6`.} \\
\texttt{`6 + 8 = 14`.} \\
\texttt{This is also correct.} \\
\texttt{Final multiplication: `14 * {-}47`.} \\
\texttt{`10 * 47 = 470`.} \\
\texttt{`4 * 47 = 188`.} \\
\texttt{`470 + 188 = 658`.} \\
\texttt{Result: `{-}658`.} \\
\\ 
\texttt{Why was this marked as a loss? Let me re{-}read the prompt's previous examples. The previous examples had specific logic errors (looping, wrong word length). Here, the math seems correct. Is there a hidden trick?} \\
\texttt{Ah, looking at the previous questions in the history:} \\
\texttt{Q1: `(({-}5 + 9 * {-}4 {-} 0) * (4 + {-}7 + 0 * {-}5))`} \\
\texttt{`9 * {-}4 = {-}36`. `{-}5 {-} 36 = {-}41`.} \\
\texttt{`4 {-} 7 = {-}3`. `0 * {-}5 = 0`. `{-}3 + 0 = {-}3`.} \\
\texttt{`{-}41 * {-}3 = 123`. Correct.} \\
\\ 
\texttt{Q2: `(({-}9 * 7 * 7 * {-}9) + (4 * {-}9 {-} 8 {-} {-}4))`} \\
\texttt{`{-}9 * 7 = {-}63`. `{-}63 * 7 = {-}441`. `{-}441 * {-}9 = 3969`.} \\
\texttt{`4 * {-}9 = {-}36`. `{-}36 {-} 8 = {-}44`. `{-}44 {-} {-}4 = {-}40`.} \\
\texttt{`3969 + {-}40 = 3929`. Correct.} \\
\\ 
\texttt{Q3: `(({-}} \\
            }
        }
    }
    & \\ \\

\end{supertabular}
}

\end{document}
